\documentclass[11pt,letterpaper]{article}
\usepackage[top=0.75in,bottom=0.75in,left=0.85in,right=0.85in]{geometry}
\usepackage{times}
\usepackage{booktabs}
\usepackage{graphicx}
\usepackage{amsmath}
\usepackage{enumitem}
\usepackage{hyperref}
\usepackage{xcolor}
\usepackage{parskip}
\usepackage{multicol}
\usepackage{caption}

\setlength{\parskip}{3.5pt}
\setlength{\parindent}{0pt}
\setlength{\abovecaptionskip}{4pt}
\setlength{\belowcaptionskip}{0pt}

\hypersetup{colorlinks=true,linkcolor=blue!60!black,urlcolor=blue!60!black,citecolor=blue!60!black}

\captionsetup{font=small}

\begin{document}

%% ── Header ───────────────────────────────────────────────────────────────────
\begin{center}
  {\large\textbf{Tree of Thoughts for Automated Code Debugging}}\\[3pt]
  {\normalsize Ronald Feng \quad Jefferson Zhou \quad Haochen Wang}\\[1pt]
  {\small Cornell University \quad CS 4782: Intro to Deep Learning \quad May 12, 2026}
\end{center}
\vspace{-4pt}\hrule\vspace{5pt}

%% ── 1. Introduction ──────────────────────────────────────────────────────────
\section*{1. Introduction}

Large language models (LLMs) perform well at many coding tasks, yet they struggle with
\emph{debugging}—a process that requires exploring competing hypotheses and backtracking from
wrong diagnoses. Standard input-output (IO) or chain-of-thought (CoT) prompting commits to a
single reasoning path, which fails when the first guess is wrong.

We re-implement and extend the \textbf{Tree of Thoughts (ToT)} framework~[1] to automated
Python code debugging. ToT organizes LLM reasoning as a tree of intermediate \emph{thoughts},
enabling systematic search (BFS or DFS) with deliberate lookahead and backtracking. Code
debugging is a natural fit: test pass/fail provides an unambiguous, automatic reward signal
that replaces the noisy LLM self-evaluation used in the original paper.
As a novel contribution, we introduce a \textbf{Monte Carlo Tree Search (MCTS)} variant with
UCB1-guided hypothesis selection and execution-based reward backpropagation.

%% ── 2. Chosen Result ─────────────────────────────────────────────────────────
\section*{2. Chosen Result}

We target Table~2 of Yao et al.~[1]: on Game of 24, \textbf{ToT-BFS (b=5) achieves 74\%} while
IO achieves only 7.3\% and CoT 4\%---an 18$\times$ gain from structured search. This result
demonstrates that deliberate tree search over intermediate reasoning steps yields dramatic
improvements over linear prompting strategies on tasks requiring multi-step planning.

We reproduce this structured-search advantage in code debugging: by exploring multiple
bug hypotheses before committing to a fix, ToT should substantially outperform single-path
baselines that must get the diagnosis right on the first attempt.

%% ── 3. Methodology ───────────────────────────────────────────────────────────
\section*{3. Methodology}

\textbf{Dataset.} We evaluate on \textbf{HumanEval-Bugs}---OpenAI's 164-problem HumanEval
benchmark with systematic single-line mutations across four categories:
\texttt{wrong\_operator} (106), \texttt{incorrect\_return} (33), \texttt{missing\_condition} (17),
\texttt{off\_by\_one} (8). Tests are the original HumanEval assertions, executed via sandboxed
\texttt{subprocess} with a 5-second timeout.

\textbf{Model and infrastructure.} All methods use \textbf{Qwen2.5-Coder-7B-Instruct} running
locally via Ollama---fully open-source, \$0 cost. Experiments run on Google Colab (NVIDIA T4 GPU).

\textbf{ToT tree structure.} Our tree has two levels (depth~=~2):
\begin{itemize}[noitemsep,topsep=1pt]
  \item \textbf{Level 1 -- Hypotheses:} The LLM generates $k{=}3$ numbered bug-diagnosis candidates
        (e.g., ``loop uses \texttt{<} instead of \texttt{<=}''). Each is scored using a
        \emph{hybrid evaluator}: 60\% LLM plausibility score + 40\% execution signal (quick-fix
        then run tests). This replaces the paper's pure LLM sure/likely/impossible scheme with
        execution feedback to reduce evaluator noise.
  \item \textbf{Level 2 -- Fixes:} For each selected hypothesis, the LLM generates $k{=}3$
        complete corrected function implementations. Each is executed against the test suite;
        the first passing fix is returned.
\end{itemize}

\textbf{Search strategies:}
\begin{itemize}[noitemsep,topsep=1pt]
  \item \textbf{BFS:} Expands all $k$ hypotheses in parallel, scores each, then generates
        and tests fixes for all survivors simultaneously. Maximum $k \times k = 9$ fix attempts.
  \item \textbf{DFS:} Greedily selects the highest-scored hypothesis, tries its $k$ fixes;
        backtracks to the next hypothesis if all fail. Identical node budget, lower token cost.
  \item \textbf{MCTS (novel):} Seeds a hypothesis pool from $k$ initial candidates, then runs
        $s{=}12$ simulations. Each simulation uses UCB1 to select a hypothesis
        ($\mathrm{UCB1}(v) = w_v/n_v + \sqrt{2}\sqrt{\ln n_{\text{parent}}/n_v}$),
        generates fixes, executes them to obtain a reward $\in\{0, 0.1, 1.0\}$, and
        backpropagates through the tree. This allocates more rollouts to promising hypotheses
        rather than exploring all branches uniformly.
\end{itemize}

\textbf{Baselines:} IO (direct fix, no reasoning), CoT (reasoning chain then fix),
CoT-SC (5 CoT samples, best fix by majority vote).

%% ── 4. Results & Analysis ────────────────────────────────────────────────────
\section*{4. Results \& Analysis}

\begin{table}[h!]
\centering
\small
\caption{All methods on HumanEval-Bugs ($n{=}164$, Qwen2.5-Coder-7B-Instruct, Google Colab T4).}
\label{tab:results}
\begin{tabular}{@{}lcccc@{}}
\toprule
\textbf{Method} & \textbf{Fix Rate} & \textbf{1st-Attempt} & \textbf{Avg Tokens} & \textbf{Avg Time} \\
\midrule
IO                         & 87.8\%           & 87.8\%               & 567                 & 2.0\,s \\
CoT                        & 77.4\%           & 77.4\%               & 940                 & 5.0\,s \\
CoT-SC ($n{=}5$)           & 95.7\%           & 95.7\%               & 4,749               & 24.6\,s \\
ToT-BFS ($k{=}3$)          & 98.8\%           & 79.9\%               & 2,874               & 14.7\,s \\
ToT-DFS ($k{=}3$)          & 98.8\%           & 86.6\%               & 2,798               & 14.2\,s \\
\textbf{MCTS} ($k{=}3$, $s{=}12$) & \textbf{100\%} & \textbf{92.1\%} & 14,335 & 64.0\,s \\
\bottomrule
\end{tabular}
\end{table}

\textbf{ToT outperforms IO and CoT.}
ToT-BFS and ToT-DFS both reach \textbf{98.8\%}, a +11 pp gain over IO.
MCTS achieves a perfect \textbf{100\%} fix rate by allocating more simulations to
the most promising hypotheses, confirmed by execution-based backpropagation.

\textbf{CoT underperforms IO (77.4\% vs.\ 87.8\%)}---a counterintuitive finding explained by
model scale. CoT was validated on GPT-4; for a 7B model on single-line bugs,
explicit reasoning chains introduce compounding errors and push the model away from its
strongest mode (direct code pattern-matching). This mirrors known results that CoT helps only
above a certain model scale threshold.

\textbf{MCTS explores 46.8 nodes on average} per problem, reallocating rollouts to promising
hypotheses rather than evaluating all branches equally. Execution-based backpropagation
corrects noisy initial LLM scores: poor hypotheses receive fewer subsequent rollouts after their
fixes fail, concentrating the search budget on viable candidates.

\textbf{DFS backtracking reveals hypothesis difficulty.}
Of 164 problems, 13 (8\%) required at least one backtrack; among those, DFS still solved
84.6\%. The 2 DFS failures exhausted all 3 hypotheses and all 9 fixes---the LLM did not
generate the correct diagnosis under any scoring order, suggesting the hybrid evaluator
occasionally misranks hypotheses for structurally ambiguous bugs.

\textbf{Limitation: dataset difficulty.}
Single-line synthetic mutations are largely solvable without structured search---IO already
reaches 87.8\%. The Tree of Thoughts advantage is most visible in token efficiency
(ToT uses 5$\times$ fewer tokens than CoT-SC for higher accuracy) and in MCTS's 100\%
ceiling. A harder dataset with multi-step logic errors (e.g., DebugBench Logic Error category)
would better surface the backtracking advantage central to the original paper.

%% ── 5. Reflections ───────────────────────────────────────────────────────────
\section*{5. Reflections}

\textbf{What worked.} Execution-based evaluation proved far more reliable than LLM
self-scoring for hypothesis ranking. The sandboxed test harness provided unambiguous rewards,
enabling MCTS backpropagation to correct initial mis-rankings over successive rollouts.
MCTS is our strongest novel result: it is the only method to achieve 100\%, and its UCB1
allocation demonstrates meaningful improvement over uniform BFS expansion.

\textbf{What we would do differently.} The core lesson is dataset selection.
Single-line mutations do not require multi-step reasoning---the tree degenerates because
the first hypothesis is almost always correct. We have updated the implementation to use
\textbf{DebugBench Logic Error} bugs (real LeetCode algorithm errors) with a \textbf{depth-3
tree} (Area $\to$ Hypothesis $\to$ Fix, matching the paper's Game of 24 structure exactly)
and \textbf{Qwen2.5-7B} (general-purpose) to provide a harder task where structured search
genuinely matters. With $k{=}5$ and depth~3, the BFS tree has up to 125 hypothesis-fix paths,
giving backtracking real opportunity to demonstrate value.

\textbf{Future directions.}
\begin{itemize}[noitemsep,topsep=1pt]
  \item Scale to repository-level bugs requiring multi-file context (SWE-bench).
  \item Incorporate dynamic analysis (stack traces, runtime values) as additional evaluation signal.
  \item Reduce MCTS token cost via early-stopping when reward exceeds threshold.
  \item Fine-tune the LLM evaluator on debugging-specific training data to improve hypothesis ranking.
\end{itemize}

%% ── 6. References ────────────────────────────────────────────────────────────
\section*{References}
\begin{enumerate}[label={[\arabic*]},noitemsep,topsep=2pt,leftmargin=*]
  \item S.\ Yao, D.\ Yu, J.\ Zhao, I.\ Shafran, T.\ L.\ Griffiths, Y.\ Cao, K.\ Narasimhan.
        ``Tree of Thoughts: Deliberate Problem Solving with Large Language Models.''
        \emph{NeurIPS 2023}.
  \item M.\ Chen et al.\ ``Evaluating Large Language Models Trained on Code (HumanEval).''
        \emph{arXiv:2107.03374}, 2021.
  \item R.\ Tian et al.\ ``DebugBench: Evaluating Debugging Capability of Large Language Models.''
        \emph{ACL 2024}.
  \item X.\ Wang et al.\ ``Self-Consistency Improves Chain of Thought Reasoning in Language Models.''
        \emph{ICLR 2023}.
  \item R.\ Coulom. ``Efficient Selectivity and Backup Operators in Monte-Carlo Tree Search.''
        \emph{CG 2006}.
  \item Qwen Team. ``Qwen2.5: A Party of Foundation Models.'' \emph{arXiv:2412.15115}, 2024.
\end{enumerate}

\end{document}
